\problemname{Sidnumrering}

\begin{center}
\includegraphics[width=0.3\textwidth]{bokbild.pdf}
\end{center}
 
I en viss bok med $N$ sidor vill förlaget spara pengar genom att enbart trycka
sidnummer på höger sida, d.v.s.\ de udda talen. Skriv ett program som
räknar ut hur många siffror som går åt av varje sort.

\section*{Indata}
En rad med heltalet $N$.

\section*{Utdata}
En rad med $10$ heltal, antalet siffror av varje sort. Först antalet nollor, sedan antalet ettor, o.s.v.

\section*{Förklaring till exempel 1}

Sidnumren som skrivs är 1, 3, 5, 7, 9, 11, 13, 15,
17, 19, 21 och 23. Alltså behövs det inga nollor, 8 ettor, 2 tvåor, 3
treor, inga fyror, 2 femmor, inga sexor, 2 sjuor, inga åttor och 2 nior.

\section*{Poängsättning}
För testfall värda $60$ poäng gäller att $1 \le N \le 100\,000$ \\
För testfall värda $40$ poäng gäller att $10^{9} \le N \le 10^{12}$ \\
